% =============================================================================================
% Seccion "Indice de Infraestructura en Argentina" del documento de INECO-UADE "Propuesta de
% Indicadores Economicos" (octubre 2026, en Overleaf; el documento completo NO esta en este repo).
% Reemplaza "\section{Índice de Infraestructura en Argentina} % Escribir aquí".
% Redactada como PROPUESTA: version preliminar del modulo nacional + indicadores y modulos pendientes.
%
% - Figuras (subirlas a la carpeta figuras/ del Overleaf; vienen en el ZIP del notebook 01):
%     figuras/iip_g02_pilares.png
%     figuras/iip_g04_telecom_pares.png
% - Paquetes: los del documento (amsmath, booktabs, graphicx, babel spanish).
% - Citas que YA estan en el documento: indec_eph, oecd2008, mecon_imig.
%   Nuevas: wef_gci, cammesa, se_gas, indec_issp, enacom_internet, ookla, bm_wdi, ocde_broadband,
%   indec_isac -> los \bibitem estan al final de este archivo.
% - Cifras: datos hasta 2026T2, commit d99d860 (verificado identico en Colab, 2026-10-08). Si se
%   actualizan los datos, revisar las cifras de ESTADO.md, seccion 2 ("donde se citan").
% =============================================================================================

\section{Índice de Infraestructura en Argentina}

La infraestructura ---energía disponible y confiable, rutas, trenes y puertos para mover la producción, conectividad, agua y cloacas--- es una condición para que las empresas produzcan, crezcan y decidan invertir, y buena parte de ella depende de la obra pública. Se propone un \emph{Índice de Infraestructura Productiva} (IIP) que resuma su estado en un valor trimestral, comparable en el tiempo, entre provincias y con otros países. Los rankings internacionales disponibles no cumplen esa función: se publican con años de rezago y el más citado, el pilar de infraestructura del Índice de Competitividad Global del Foro Económico Mundial, no se actualiza desde 2019 \cite{wef_gci}. A diferencia de los demás indicadores de este documento, el IIP está en una etapa inicial: lo que sigue es una propuesta metodológica acompañada de una primera versión calculada a nivel nacional, cuyos indicadores, pesos y alcance deberían revisarse con especialistas de cada sector antes de su publicación.

\subsection{Propuesta metodológica}

El índice se organiza en cinco pilares (Cuadro~\ref{tab:iip_pilares}). Los cuatro primeros miden el estado de la infraestructura ---su cobertura, capacidad, calidad y uso--- y el quinto, la inversión que la repone y la amplía, que anticipa la evolución de los demás. La versión preliminar usa doce indicadores trimestrales de fuentes públicas desde 2016; la última columna del cuadro reúne candidatos a incorporar, algunos de los cuales todavía no cuentan con una fuente abierta y actualizada.

\begin{table}[!htb]
\centering
\small
\caption{Pilares del IIP: indicadores de la versión preliminar y candidatos a incorporar}
\vspace{.25 cm}
\label{tab:iip_pilares}
\begin{tabular}{@{}p{3.0cm}p{6.3cm}p{5.8cm}@{}}
\toprule
\textbf{Pilar} & \textbf{Versión preliminar} & \textbf{A incorporar} \\
\midrule
Energía & Margen de reserva eléctrica: potencia instalada sobre el pico de demanda de 36 meses \cite{cammesa}; autoabastecimiento de gas: producción sobre consumo \cite{se_gas} & Calidad del servicio eléctrico (duración y frecuencia de los cortes); hogares con gas de red \\
Transporte y logística & Carga ferroviaria y carga aérea por unidad de actividad \cite{indec_issp} & Estado de la red vial; movimiento de contenedores en puertos; desempeño logístico (LPI) \\
Telecomunicaciones & Accesos a internet fijo por habitante, velocidad medida de bajada y proporción de accesos por fibra óptica, frente a los países de comparación \cite{enacom_internet,ookla,bm_wdi,ocde_broadband} & Cobertura de las redes móviles 4G y 5G \\
Agua y saneamiento & Hogares con agua de red y con cloacas en los 31 aglomerados urbanos \cite{indec_eph} & Cobertura total y por provincia (Censo 2022) \\
Inversión & Gasto de capital nacional en energía, transporte y agua (\% del PIB) \cite{mecon_imig}; insumos de la obra vial (asfalto) y despachos de cemento \cite{indec_isac} & Obra pública por provincia; inversión privada en infraestructura; stock de capital público \\
\bottomrule
\end{tabular}
\end{table}

Tres decisiones guían la construcción. Primero, los indicadores de uso, como la carga transportada, siguen al ciclo económico más que a la infraestructura, por lo que se expresan por unidad de actividad (divididos por el EMAE). Segundo, en telecomunicaciones la tecnología mejora en todo el mundo: la velocidad contratada en la Argentina pasó de 3,6 a 265~Mbps entre 2014 y 2026 \cite{enacom_internet} y, medida contra su propia historia, mostraría una mejora permanente. Por eso esos indicadores se miden como brecha frente a la mediana de diez países de comparación (Brasil, Chile, Uruguay, México, Colombia, Perú, Australia, Canadá, España y Estados Unidos),
\begin{equation*}
b_t = \frac{x^{\text{ARG}}_t}{\operatorname{mediana}_i\, x^{i}_t} - 1,
\end{equation*}
con la velocidad medida por Ookla \cite{ookla} en lugar de la contratada, que no es comparable entre países. Tercero, la inversión se mide en relación con el PIB y se agrupa en un pilar propio.

La normalización es la del Índice de Fortaleza Macroeconómica Argentina, para que ambos índices se lean igual: cada indicador se convierte en un puntaje robusto con la mediana y el RIC de su historia desde 2016, recortado a $\pm 3$; cada pilar promedia sus indicadores y el índice promedia los pilares con igual peso \cite{oecd2008}. Los pesos son un parámetro del cálculo y pueden ajustarse si la evaluación de especialistas lo aconseja. Como algunos indicadores comienzan más tarde (agua y cloacas en 2017, la velocidad medida en 2019), el índice se encadena hacia atrás: el último trimestre usa todos los indicadores y la historia se mueve solo con las variaciones de los disponibles, de modo que la incorporación de un indicador no produce saltos. Las fuentes internacionales publican con rezago (la ITU y la OCDE llegan a fines de 2024), por lo que la referencia de los pares se extiende con la tendencia de cada país durante a lo sumo seis trimestres. Al igual que en el índice macroeconómico, se informa un rango de sensibilidad con trece variantes metodológicas: la exclusión de cada pilar, la comparación con el mejor de los pares en lugar de la mediana, otra ventana de normalización y otro recorte, entre otras.

La propuesta contempla tres módulos con los mismos pilares: el nacional, ya calculado; uno provincial, que ordene a las 24 jurisdicciones según su infraestructura para orientar decisiones de localización, y uno internacional, que ubique a la Argentina frente a sus pares en cada pilar. Para ambos ya se relevaron fuentes: ENACOM y Ookla publican datos de internet por provincia; el Censo 2022, la cobertura de agua, cloacas y gas de red; CAMMESA, la demanda eléctrica por provincia; el Mapa de Inversiones de la Secretaría de Obras Públicas, la obra pública por provincia, y el Banco Mundial y la OCDE, indicadores comparables entre países.

\subsection{Resultados preliminares}

En el segundo trimestre de 2026 el IIP se ubicó en $+0{,}32$ desvíos: por encima de lo típico del período en las trece variantes (rango de $+0{,}09$ a $+0{,}74$). Mejoró $0{,}18$ en el último año, también en todas las variantes (rango de $+0{,}11$ a $+0{,}27$). La mejora, sin embargo, es desigual entre pilares (Figura~\ref{fig:iip_pilares}).

\begin{figure}[!htb]
\centering
\includegraphics[width=\textwidth]{figuras/iip_g02_pilares.png}
\caption{IIP y sus pilares, 2016T1--2026T2, en desvíos respecto de lo típico del período. Las bandas grises marcan las gestiones de gobierno.}
\label{fig:iip_pilares}
\end{figure}

La energía ($+0{,}91$) y las telecomunicaciones ($+1{,}17$) explican la mejora. La producción de gas cubre el 122\% del consumo interno, frente al 99\% de 2016, por el desarrollo de Vaca Muerta. El margen de reserva eléctrica, en cambio, es del 46\%, por debajo del máximo de 62\% de 2021, porque el pico de demanda creció más que la potencia instalada. En telecomunicaciones, la Argentina alcanzó a sus pares en accesos a internet fijo (29,6 cada 100 habitantes, frente a 29,7 de la mediana; en 2016 tenía un 25\% menos) y acortó la brecha en fibra óptica: el 50\% de sus accesos, un 34\% menos que los pares, frente a un 73\% menos en 2016. La velocidad medida, en cambio, es de 180~Mbps, un 37\% por debajo de la mediana de los pares (286~Mbps), y esa brecha se amplió desde cerca del 20\% en 2019 (Figura~\ref{fig:iip_telecom}): el país está más conectado que antes, pero con redes relativamente más lentas.

\begin{figure}[!htb]
\centering
\includegraphics[width=\textwidth]{figuras/iip_g04_telecom_pares.png}
\caption{Telecomunicaciones: Argentina frente a la mediana de los diez países de comparación. El tramo punteado extiende la tendencia de los pares mientras las fuentes internacionales no publican.}
\label{fig:iip_telecom}
\end{figure}

En agua y saneamiento ($+0{,}87$), el 91,0\% de los hogares urbanos tiene agua de red y el 73,0\%, cloacas; los avances son de menos de un punto, pero la estandarización los amplifica porque estas series se mueven poco, lo que justifica sumar el Censo 2022 para medir la cobertura total. El transporte ($+0{,}02$) mueve una carga típica por unidad de actividad. El pilar de inversión ($-1{,}37$) está en el mínimo de la serie: el gasto de capital nacional en energía, transporte y agua fue de 0,19\% del PIB en los doce meses a junio de 2026, frente a 0,80\% en 2023 y 1,06\% en 2016, y los insumos de la obra vial y los despachos de cemento están por debajo de lo típico.

\begin{table}[!htb]
\centering
\small
\caption{IIP y pilares: promedio por gestión (desvíos respecto de lo típico)}
\vspace{.25 cm}
\label{tab:iip_gestiones}
\begin{tabular}{@{}l r r r r r r@{}}
\toprule
Gestión & IIP & Energía & Transporte & Telecom. & Agua & Inversión \\
\midrule
Macri (2016T1--2019T4)        & $-0{,}19$ & $-0{,}89$ & $+0{,}06$ & $+0{,}02$ & $-0{,}55$ & $+0{,}72$ \\
A.~Fernández (2020T1--2023T4) & $+0{,}11$ & $+0{,}56$ & $-0{,}08$ & $+0{,}28$ & $-0{,}28$ & $+0{,}06$ \\
Milei (2024T1--2026T2)        & $+0{,}14$ & $+0{,}58$ & $-0{,}12$ & $+0{,}74$ & $+0{,}70$ & $-1{,}18$ \\
\bottomrule
\end{tabular}

\smallskip
\parbox{13.5cm}{\footnotesize Promedios de los valores trimestrales. Agua y saneamiento desde 2017T4. Fuente: INECO--UADE en base a CAMMESA, Secretaría de Energía, ENARGAS, INDEC, ENACOM, Ookla, Banco Mundial (ITU), OCDE y Secretaría de Hacienda.}
\end{table}

Promediado por gestión, el índice da $-0{,}19$ (Macri), $+0{,}11$ (A.~Fernández) y $+0{,}14$ (Milei, hasta el segundo trimestre de 2026), como muestra el Cuadro~\ref{tab:iip_gestiones}. El promedio de Macri es el más bajo en las trece variantes, pero la infraestructura es un stock: ese promedio incluye el punto de partida de 2016, con un margen de reserva del 33\% y un déficit de gas, y sin el pilar de energía la diferencia es mínima. El orden entre A.~Fernández y Milei no es robusto. Como en el índice macroeconómico, estos promedios describen condiciones y no evalúan gestiones.

En síntesis, los resultados preliminares indican que la infraestructura productiva está algo mejor que lo típico de la última década y que mejoró en el último año, sostenida por sectores donde predomina la inversión privada, como el gas y las telecomunicaciones. Al mismo tiempo, la inversión pública en infraestructura está en mínimos. Como la infraestructura se deteriora sin mantenimiento, el pilar de inversión anticipa presiones futuras sobre el transporte y el agua, donde el Estado es el principal inversor.

\subsection{Próximos pasos}

Para pasar de la propuesta a un indicador publicable se sugiere:
\begin{itemize}
\item Revisar los indicadores, los pesos y los países de comparación con especialistas en energía, transporte, telecomunicaciones y agua, en el marco de las comisiones de evaluación propuestas en la Introducción.
\item Incorporar los indicadores pendientes del Cuadro~\ref{tab:iip_pilares}, en especial la calidad del servicio eléctrico y el estado de la red vial, para los que habría que gestionar el acceso a los datos con el ENRE y Vialidad Nacional.
\item Desarrollar el módulo provincial, con un ranking de las 24 jurisdicciones, y el internacional, con la posición de la Argentina frente a sus pares en cada pilar.
\item Validar el índice contra la inversión bruta interna fija y la actividad de los sectores más intensivos en infraestructura, y publicarlo trimestralmente en el boletín de INECO.
\end{itemize}

% =============================================================================================
% \bibitem nuevos (pegar dentro de \begin{thebibliography} ... \end{thebibliography}).
% indec_eph, oecd2008 y mecon_imig ya estan en el documento.
%
% \bibitem{wef_gci} World Economic Forum (2019). \textit{The Global Competitiveness Report 2019}. Ginebra: WEF.
%
% \bibitem{cammesa} CAMMESA (2026). \textit{Demanda de energía eléctrica y potencia instalada del Mercado Eléctrico Mayorista, series mensuales}. Vía Subsecretaría de Programación Macroeconómica, datos.gob.ar.
%
% \bibitem{se_gas} Secretaría de Energía y ENARGAS (2026). \textit{Producción y consumo de gas natural, series mensuales}. Vía Subsecretaría de Programación Macroeconómica, datos.gob.ar.
%
% \bibitem{indec_issp} INDEC (2026). \textit{Estadísticas de servicios públicos}. Buenos Aires: Instituto Nacional de Estadística y Censos.
%
% \bibitem{enacom_internet} ENACOM (2026). \textit{Indicadores de internet fijo: accesos, penetración, velocidad y tecnologías}. Ente Nacional de Comunicaciones, indicadores.enacom.gob.ar.
%
% \bibitem{ookla} Ookla (2026). \textit{Speedtest\textsuperscript{\textregistered} by Ookla\textsuperscript{\textregistered} Global Fixed and Mobile Network Performance Maps}. Datos abiertos, licencia CC BY-NC-SA 4.0.
%
% \bibitem{bm_wdi} Banco Mundial (2026). \textit{World Development Indicators} (suscripciones de banda ancha fija con datos de la Unión Internacional de Telecomunicaciones). Washington, DC: Banco Mundial.
%
% \bibitem{ocde_broadband} OCDE (2026). \textit{Broadband and Telecom Database}. París: Organización para la Cooperación y el Desarrollo Económicos.
%
% \bibitem{indec_isac} INDEC (2026). \textit{Indicador sintético de la actividad de la construcción (ISAC)}. Buenos Aires: Instituto Nacional de Estadística y Censos.
% =============================================================================================
